% Part: many-valued-logic
% Chapter: syntax-and-semantics
% Section: formulas

\documentclass[../../../include/open-logic-section]{subfiles}

\begin{document}

\olfileid[id]{mvl}{syn}{fml}

\olsection{\usetoken{P}{formula}}

\begin{defn}[Formula]
\ollabel{defn:formulas}
Himpunan~$\Frm[L]$ yang beranggotakan \emph{!!{formula}s} dari suatu bahasa
proposisional~$\Lang L$ didefinisikan secara induktif sebagai berikut:
\begin{enumerate}
\item Setiap !!{propositional variable}~$\Obj p_i$ merupakan !!{formula}
  atomik.
\item Setiap konektif $0$-tempat (konstanta proposisional) dari~$\Lang L$
  merupakan !!{formula} atomik.
\item Jika $\star$ merupakan konektif $n$-tempat dari~$\Lang L$, dan $!A_1$,
  \dots, $!A_n$ merupakan !!{formula}s, maka $\star(!A_1,\dots,!A_n)$
  merupakan !!a{formula}.
\tagitem{limitClause}{Tidak ada yang lain yang merupakan !!a{formula}.}{}
\end{enumerate}
Jika $\star$ merupakan konektif $1$-tempat, maka $\star(!A_1)$ sering ditulis
sebagai $\star !A_1$. Jika $\star$ merupakan konektif $2$-tempat, maka
$\star(!A_1,!A_2)$ sering ditulis sebagai $(!A_1\star !A_2)$.
\end{defn}

Seperti biasa, kita sering menghilangkan secara diam-diam pasangan kurung
terluar.

\begin{ex}
  Dalam bahasa baku~$\Lang{L_0}$, $\Obj p_1\lif(\Obj p_1\land\lnot\Obj p_2)$
  merupakan suatu formula. Dalam bahasa logika produk, formula itu ditulis
  sebagai $\Obj p_1\lif(\Obj p_1\odot\lnot\Obj p_2)$. Jika kita menambahkan
  konektif $1$-tempat~$\triangle$ ke dalam bahasa tersebut, kita juga akan
  memperoleh formula seperti $\triangle(\Obj p_1\land\Obj p_2)\lif
  (\triangle\Obj p_1\land\triangle\Obj p_2)$.
\end{ex}

\end{document}
